\section{Research Questions and Study Design}
\label{sec:study_design}

This section states the five research questions, the hypotheses they are
designed to test, the variables involved, and the statistical analysis plan.
The corresponding experimental setup---datasets, strategies, learners, sample
sizes, and protocol---is detailed in \Cref{sec:experiment}.

The overarching research goal of the study is:
\begin{quote}
\emph{To empirically characterize when training sample size $n$ is a valid
proxy for empirical acquisition cost, when it is miscalibrated, and whether
replacing count-based cost with empirical cost changes conclusions about
sampling strategies.}
\end{quote}
We decompose this goal into five questions, RQ1--RQ5, which together cover
heterogeneity, correlation/calibration, ranking robustness, Pareto efficiency,
and context sensitivity.

\subsection{RQ1 -- Empirical Cost Heterogeneity}
\label{subsec:rq1}

\textbf{RQ1.} \emph{At fixed training sample size $n$, how much does
empirical acquisition cost vary across configurations and across sampling
strategies?}

\paragraph*{Motivation.}
Prior work assumes, implicitly or explicitly, that all configurations have
equal acquisition cost (\Cref{subsec:cost-models}). RQ1 tests this by
holding $n$ fixed and comparing $\TrainCost(r,s,n)$ across strategies. If
at fixed $n$ some strategies systematically incur more (or less)
acquisition cost than others, the size-based cost abstraction hides a
strategy-induced bias.

\paragraph*{Hypotheses.}
\begin{itemize}
    \item \textbf{H1a.} Per-configuration acquisition cost
    $\mathit{cost}(c)$ is heterogeneous within a system, \ie its
    coefficient of variation is non-negligible.
    \item \textbf{H1b.} At fixed $n$, the distribution of $\TrainCost$ differs
    across sampling strategies.
    \item \textbf{Null.} At fixed $n$, sampling strategy has no significant
    effect on $\TrainCost$.
\end{itemize}

\paragraph*{Variables.}
Independent: system, input/workload (where applicable), property $p$, strategy
$s$, training size $n$, run $r$. Dependent: $\TrainCost(r,s,n)$. Secondary
descriptors: per-configuration cost distribution, coefficient of variation,
overhead relative to random.

\paragraph*{Analysis.}
For each (system, input, property, $n$), we
\begin{enumerate}
    \item report the median and IQR of $\TrainCost$ per strategy across runs;
    \item compare strategies using the Kruskal--Wallis
    test~\cite{kruskal1952};
    \item perform pairwise Mann--Whitney $U$ tests~\cite{mann1947} of each
    strategy against random;
    \item report Vargha--Delaney $\hat A_{12}$ effect sizes~\cite{vargha2000}
    with the usual thresholds (small$>0.56$, medium$>0.64$, large$>0.71$);
    \item report a practical \emph{overhead} relative to random,
    $\mathit{Overhead}(s,n)=\mathrm{median}\,\TrainCost(s,n)/\mathrm{median}\,\TrainCost(\mathrm{Random},n)$,
    classified as negligible ($<5\%$), small ($5$--$15\%$), moderate
    ($15$--$30\%$), or substantial ($>30\%$).
\end{enumerate}

\paragraph*{Interpretation.}
If strategies differ significantly at fixed $n$, then ``cost = sample size''
hides a strategy-induced cost bias. If they do not, $n$ dominates cost in that
regime.

\subsection{RQ2 -- Validity of Sample Size as a Cost Proxy}
\label{subsec:rq2}

\textbf{RQ2.} \emph{To what extent is $n$ a valid proxy for $\TrainCost$, in
terms of correlation, calibration, and strategy-specific bias?}

\paragraph*{Motivation.}
Even if cost grows with $n$, the relationship may not be calibrated: a
strategy spending $20\%$ of $n$ might consume $10\%$, $20\%$, or $40\%$ of
the full acquisition cost. RQ2 distinguishes between rank correlation and
budget calibration.

\paragraph*{Hypotheses.}
\begin{itemize}
    \item \textbf{H2a.} $n$ is positively rank-correlated with $\TrainCost$.
    \item \textbf{H2b.} The calibration ratio
    $\Ratio = \mathit{TrainCostPct}/\mathit{SizePct}$ is not centered at $1$.
    \item \textbf{H2c.} Calibration depends on system, strategy, input and
    property.
\end{itemize}

\paragraph*{Metrics.}
Spearman correlation $\rho(n,\TrainCost)$ per (system, input, property,
strategy); the calibration ratio of \Cref{eq:ratio}; the strategy-specific bias
$\Bias(s)$ of \Cref{eq:bias}; and the slope/intercept of an OLS regression of
$\mathit{TrainCostPct}$ on $\mathit{SizePct}$.

\paragraph*{Analysis.}
We report the joint distribution of $(\mathit{SizePct},\mathit{TrainCostPct})$
as scatter plots with the $y=x$ identity line, fit a linear/log-linear
regression, and report median and IQR of the calibration ratio per
strategy. Pairwise $\Bias$ comparisons across strategies use Mann--Whitney $U$.

\paragraph*{Interpretation.}
We expect four qualitative outcomes:
\begin{enumerate}
    \item high correlation and ratio close to $1$: $n$ is a good proxy;
    \item high correlation but ratio far from $1$: $n$ ranks costs but
    miscalibrates budgets;
    \item low/moderate correlation: $n$ is unreliable as a cost proxy;
    \item strategy-specific bias: some strategies systematically select cheaper
    or costlier configurations.
\end{enumerate}

\subsection{RQ3 -- Robustness of Conclusions to the Cost Model}
\label{subsec:rq3}

\textbf{RQ3.} \emph{Do strategy rankings under accuracy--cost trade-offs change
when empirical training cost replaces count-based cost?}

\paragraph*{Motivation.}
Even if $n$ is correlated and approximately calibrated with $\TrainCost$ in
aggregate, conclusions may still flip when strategies are ranked by
$J^{\text{size}}_\lambda$ vs.\ $J^{\text{time}}_\lambda$, because the ranking
depends on relative ordering rather than on aggregate calibration. RQ3 directly
tests the scientific consequence of the size-as-cost abstraction: does it
ever change which strategy you would choose?

\paragraph*{Hypotheses.}
\begin{itemize}
    \item \textbf{H3a.} Rankings under count-based and empirical-cost
    objectives sometimes diverge.
    \item \textbf{H3b.} Divergence is stronger when per-configuration cost
    heterogeneity is high or when strategies select cost-biased regions of the
    configuration space.
    \item \textbf{Null.} Rankings remain stable across the two cost models.
\end{itemize}

\paragraph*{Analysis.}
For each (system, input, property), and for each $\lambda$ in the sweep
$\{0.1,0.25,0.5,1,2,4,10\}$, we
\begin{enumerate}
    \item rank strategies (or strategy/$n$ pairs) by $J^{\text{size}}_\lambda$
    and by $J^{\text{time}}_\lambda$;
    \item compute Kendall's $\tau$ between the two rankings;
    \item compute top-$1$ and top-$3$ agreement;
    \item compute the mean rank difference and identify rank flips.
\end{enumerate}

\paragraph*{Interpretation.}
A small mean rank difference and high $\tau$ across $\lambda$ values mean
size-based conclusions are robust in that context. Rank flips, especially at
the top, are evidence that the cost-model choice matters for the
recommendation.

\subsection{RQ4 -- Pareto Efficiency under Empirical Cost}
\label{subsec:rq4}

\textbf{RQ4.} \emph{Which $(\text{strategy},n)$ pairs are Pareto-efficient under
empirical training cost, and how different is the empirical-cost frontier from
the count-based one?}

\paragraph*{Motivation.}
Composite objectives require choosing $\lambda$. A complementary,
parameter-free way to ask whether the cost model matters is to compare
Pareto frontiers in the MRE--cost plane.

\paragraph*{Hypotheses.}
\begin{itemize}
    \item \textbf{H4a.} Some $(s,n)$ pairs are Pareto-efficient under
    count-based cost but dominated under empirical cost (or vice versa).
    \item \textbf{H4b.} The empirical-cost frontier differs from the
    count-based frontier in at least some systems/contexts.
    \item \textbf{Null.} The two frontiers coincide or differ only marginally.
\end{itemize}

\paragraph*{Analysis.}
For each (system, input, property) we compute and compare two frontiers:
$\{(n,\MRE)\}$ and $\{(\TrainCost,\MRE)\}$. We report the number of
Pareto-efficient pairs on each frontier, the membership overlap, the pairs that
appear only on one frontier, and the resulting frontier expansion or
contraction under measured cost. Following Li \etal's guidance that Pareto
evaluation in SBSE should match the concrete decision question, we do not lead
with generic approximation indicators here: the task is to determine whether
the same finite set of $(s,n)$ candidates remains non-dominated when the cost
axis changes \cite{li2020pareto}. This is the most decision-oriented of our
analyses: it tells practitioners which $(s,n)$ choices are not obviously
wasteful, regardless of $\lambda$.

\subsection{RQ5 -- Context Sensitivity of Cost-Model Validity}
\label{subsec:rq5}

\textbf{RQ5.} \emph{When do the effects characterized in RQ1--RQ4 appear or
disappear across systems, inputs and non-functional properties?}

\paragraph*{Motivation.}
Context sensitivity of sampling-based performance prediction is itself
established prior art (\Cref{subsec:context-sensitivity}). The novel angle is
the context sensitivity of \emph{cost-model validity}. The x264 case study is
ideal here: x264 is the most-used benchmark in this literature, has $1{,}152$
valid configurations and $16$ binary options, and admits a natural design over
$17$ input videos and two performance properties (encoding time and output
size). The asymmetric case (predicting size but paying time-based cost) is
particularly informative because the predicted property and the cost property
are not the same quantity.

\paragraph*{Hypotheses.}
\begin{itemize}
    \item \textbf{H5a.} The validity of $n$ as a cost proxy depends on the
    input/workload.
    \item \textbf{H5b.} The size-as-cost gap differs between predicting time
    and predicting size on x264.
    \item \textbf{H5c.} Strategy rankings and Pareto frontiers are less stable
    under empirical cost than under count-based cost.
\end{itemize}

\paragraph*{Analysis.}
For each input and property we repeat RQ1--RQ4. We summarize cross-context
stability via Kendall's $W$ over rankings, the average Kendall's $\tau$ across
context pairs, and the frontier overlap across contexts. We then cluster
contexts into ones where size-based cost is safe and ones where it is not, and
discuss what distinguishes them (\eg cost coefficient of variation per input,
correlation between option activation and cost).

\subsection{Expected Result Patterns}

To guide the interpretation in \Cref{sec:results,sec:discussion}, we
list the qualitative patterns we expect and how each would be read.
\begin{itemize}
    \item \emph{Pattern A.} $n$ valid everywhere: high correlation, ratios near
    $1$, no rank flips, identical frontiers. The size-as-cost abstraction is
    empirically justified for these systems. This is still a contribution
    because the abstraction had not been tested.
    \item \emph{Pattern B.} Correlated but miscalibrated: $\rho$ high but
    ratios away from $1$. Size ranks cost approximately but misestimates
    budgets.
    \item \emph{Pattern C.} Strategy-induced bias: at fixed $n$, strategies
    differ significantly in $\TrainCost$. Which configurations are selected
    matters, not only how many.
    \item \emph{Pattern D.} Rankings or frontiers change: empirical cost alters
    the best strategy or the Pareto-efficient set. Size-based conclusions are
    sensitive to the cost abstraction.
    \item \emph{Pattern E.} Effects appear only for some x264 inputs or
    properties. Cost-model validity is context-dependent and x264 should not
    be treated as a single universal benchmark point.
\end{itemize}
None of these patterns is incompatible with the others; the results may show
mixtures across systems, inputs and properties, and that mixture is itself
part of the answer.

\section{Experimental Setup}
\label{sec:experiment}

This section details how RQ1--RQ5 are operationalized: datasets
(\Cref{subsec:datasets}), strategies and baselines (\Cref{subsec:strategies}),
sample sizes (\Cref{subsec:sizes}), learners (\Cref{subsec:learners}), and the
evaluation protocol (\Cref{subsec:protocol}).

\subsection{Datasets}
\label{subsec:datasets}

A central design constraint of this study is that empirical
$\TrainCost$ requires per-configuration acquisition costs with a
documented measurement protocol (warm-up, repetition, aggregation). In
this paper we instantiate $\mathit{cost}(c)$ as wall-clock measurement
time, the most commonly constrained resource in practice; the
formulation supports other quantitative resources (\Cref{subsec:empirical-cost}).
We prioritize datasets that satisfy this constraint and structure them in
three tiers.

\paragraph*{Tier 1A: x264 with $17$ inputs and $2$ properties (ICPE 2020 artifact).}
x264 is the canonical benchmark of this literature and the most-used case
study in the systematic review of Pereira \etal~\cite{julianasurvey}. We
adopt the configuration table of Kaltenecker
\etal~\cite{kaltenecker2019}: $1{,}152$ valid configurations over $16$
binary options. Following Pereira \etal~\cite{pereira2020}, we measure each
configuration on $17$ input videos covering a range of resolutions,
motion content and durations (the full list of inputs and their
measurement-stability values is in \Cref{app:input-videos}). We consider
two non-functional properties: \emph{encoding time} (in seconds) and
\emph{output size} (in bytes of the H.264 output). For each (configuration,
input, property), we record the property value as well as the elapsed
encoding time used as $\mathit{cost}(c)$. When the predicted
property is \emph{time}, the target and the cost are the same physical
quantity. When the predicted property is \emph{size}, the target and the
cost are different quantities, because measuring output size still requires
running the encoder---this is the asymmetric case discussed in
\Cref{sec:cost_model} and \Cref{subsec:rq5}.

\paragraph*{Tier 1B: x264 on the YouTube UGC corpus.}
To stress-test input sensitivity at a much larger scale, we additionally
use the public x264-on-YouTube-UGC dataset
\red{[cite Lesoil et al., Zenodo 3928253; bib entry needed]}: $201$
configurations measured on $1{,}397$ user-generated videos. If sample count
fails as a cost proxy anywhere, this is where it should fail most
visibly. We use this corpus for RQ5 only and report it separately from
Tier-1A so that detailed comparisons remain reproducible against
ICPE\,2020.

\paragraph*{Tier 2A: nine additional configurable systems (Kaltenecker corpus).}
We extend the analysis to nine other configurable systems used in
Kaltenecker \etal~\cite{kaltenecker2019} and originating from
Siegmund \etal{} and others~\cite{siegmund2015FSE,julianasurvey} for which
both a feature model and per-configuration measured times are publicly
available. Seven of these systems are reported in our cross-system
results (\Cref{tab:systems-and-cost}, \Cref{sec:results}), spanning
$|C|\in[432,\;68{,}640]$ valid configurations and $|O|\in[11,\;54]$
options. JavaGC and VP9 are not yet included in this batch.

\paragraph*{Tier 2B: \dsWLI ICSE 2023 package.}
For workload sensitivity, we additionally use the \dsWLI dataset
of \red{[cite Mühlbauer et al., ICSE 2023; bib entry needed]}, which
provides $25{,}258$ measurements across nine configurable systems
(jump3r, kanzi, dconvert, h2, batik, xz, lrzip, x264, z3) on $6$--$13$
workloads each. Crucially, the package documents a five-run median
protocol. The H2 case is throughput-based, so we keep it in the study but
interpret the measured property carefully: when throughput is the
prediction target, we predict raw throughput; when throughput is the cost
axis, we orient it as an inverse cost so that lower values still mean
cheaper acquisition. When raw repetition logs are not exposed, we
approximate acquisition cost from the published median measurement and
explicitly label the result as an approximation in our reporting (see
\Cref{sec:threats}).

\paragraph*{Tier 2C: input-aware performance models.}
The input-aware performance models repository
\red{[cite Lesoil et al., JSS 2023/2024; bib entry needed]} provides
$1{,}941{,}075$ configuration measurements across eight systems (gcc,
ImageMagick, lingeling, nodeJS, poppler, SQLite, x264, xz) with multiple
input/workload variants per system and several measured properties
(compile time, execution time, binary size, query times, encoding
time/fps/bitrate/size). It is the largest configurable-systems collection
we use and is well-suited to RQ5 analyses that span both workload and
target property. Where a property is not time-like (e.g., binary size,
bitrate), we use the published elapsed-time column as
$\mathit{cost}(c)$ and keep workload identity attached, so that
$\TrainCost$ is computed as $\sum_{c\in S}\mathrm{time}(c, \mathrm{input})$
rather than $\sum_{c\in S}\mathrm{time}(c)$.

\paragraph*{Tier 3: robustness add-ons.}
We use two additional datasets to test the robustness of our conclusions.
\emph{Performance Evolution}
\red{[cite Kaltenecker et al., performance-evolution; bib entry needed]}
provides per-release measurements across $190$ releases of $12$ systems
(brotli, Fast Downward, HSQLDB, lrzip, MariaDB, MySQL, OpenVPN, Opus,
PostgreSQL, VP8, VP9, z3): if our conclusions hold across software
versions, that is evidence that they are not artefacts of a single
release. \emph{Green Configuration / the performance--energy corpus}
\red{[cite Weber et al.; bib entry needed]} provides $14$ systems with
runtime and energy measurements, repeated three times and averaged after
bias pre-experiments; this dataset is particularly useful for showing that
a global runtime--energy correlation can break locally, by analogy to our
size-as-cost validity question.

\subsubsection{Linux kernel cross-version study (separate protocol)}
\label{subsubsec:linux-protocol}

We additionally run a Linux kernel study on two \dsTux releases.
This study is reported \emph{separately} from \StudyReplication
throughout the paper because its protocol is structurally different
from the synthesis-based protocol of Sections~\ref{subsec:strategies}
and~\ref{subsec:learners}. Aggregating the two would conflate two
different sampling mechanisms.

\paragraph*{Pool and cost.}
We use the \dsTux artifacts published by Borges \etal\ at
EASE\,2025 \red{[bib entry needed]}: \emph{Linux v4.13} (OpenML
\texttt{46759}, $92{,}562$ valid configurations) and \emph{Linux v5.08}
(OpenML \texttt{46744}, $21{,}923$ valid configurations). Both pools
are drawn with \texttt{make randconfig} on their respective kernel
trees; the smaller v5.08 pool is a property of the published artifact,
not of the kernel itself. The cost construct here is
storage-bound: $\mathit{cost}(c)$ is the uncompressed kernel binary
size (\texttt{vmlinux}, in bytes), and the prediction target is the
same quantity. This is the only place in the paper where target $=$
cost \emph{and} both are non-time-valued; the published artifact does
not ship build time, so a time-based instance of $\mathit{cost}(c)$
on Linux is out of reach for this paper.

\paragraph*{Pool-based sampling, not synthesis-based.}
The crucial protocol difference is that Linux strategies are
\emph{pool-based subsamplers}: each strategy selects $n$ rows from a
fixed randconfig-induced pool. They do not synthesize configurations
from a feature model. The seven strategies are five deployable
samplers and two diagnostic oracles:
\begin{itemize}
    \item \texttt{random}: uniform pool subsampling (reference baseline).
    \item \texttt{balanced}: stratified subsampling on
    active-options decile.
    \item \texttt{coverage}: greedy option-frequency matching.
    \item \texttt{uncertainty}: residual-driven active learning
    (ridge model fit on a seed of $\max(50, \lceil n/4\rceil)$ rows,
    then top-$k$ residual selection).
    \item \texttt{diversity}: farthest-first in Hamming distance,
    capped at $n\le 1{,}851$ because the $O(n^{2})$ cost becomes
    prohibitive at larger sizes (we report only the $n$ levels where
    \texttt{diversity} runs).
    \item \texttt{cheapOracle}, \texttt{expensiveOracle}: lower- and
    upper-bound diagnostics that read the cost signal.
\end{itemize}
None of the synthesis-based strategies of
\Cref{subsec:strategies} (Henard, t-wise, distance-based,
solver-based, diversified distance-based) is evaluable on this design
because there is no formal model handle to drive them; conversely,
\texttt{uncertainty} and \texttt{diversity} have no analogues in the
\StudyReplication protocol. We therefore never combine the two.

\paragraph*{Sample sizes, learner, evaluation.}
Five training fractions $\{1\%, 2\%, 5\%, 10\%, 20\%\}$ of pool size,
giving $n\in\{926, 1851, 4628, 9256, 18512\}$ on v4.13 and
$n\in\{219, 438, 1096, 2192, 4385\}$ on v5.08. Stochastic strategies
are repeated $100$ times (\texttt{random}, \texttt{balanced}) or $20$
times (\texttt{uncertainty}, \texttt{diversity}, \texttt{coverage}).
The learner is the same lightweight ridge log-linear model with
$\alpha=1$ used for \StudyReplication; MRE is computed on the
complement of the training sample within the pool, with chunked
predictions ($4{,}096$ rows per chunk).

\paragraph*{Forbidden columns.}
$21$ columns are excluded from the design matrix: the target
\texttt{vmlinux}; eighteen compressed-image variants
($\{\texttt{GZIP},\texttt{BZIP2},\texttt{LZMA},\texttt{XZ},\texttt{LZO},\texttt{LZ4}\}\times\{\texttt{bzImage},\texttt{vmlinux},\texttt{raw}\}$)
which are tightly correlated with the target; \texttt{active\_options},
the count of enabled options (exactly derivable from the option
vector); and the row identifier \texttt{cid}. After near-constant
pruning (minor-class $\geq 50$) and a top-$K=512$ marginal-effect
filter, $495$ kernel options enter the v4.13 design matrix and $494$
enter v5.08.

\paragraph*{Why two releases.}
We use two releases (v4.13--v5.08; seven releases) to test whether the
calibration findings replicate across pool sizes and across kernel
generations, and whether any seams are visible across pools that
would not surface from a single release. Both questions admit
non-trivial answers (\Cref{sec:results}).

\subsubsection{the performance--energy corpus generated-sampling study (separate protocol)}
\label{subsubsec:twins-protocol}

We additionally run a study on the Green Configuration /
\emph{the performance--energy corpus} dataset \red{[bib entry needed:
Weber et al., ICSE 2023; Zenodo \texttt{7544891}]}. As with Linux, this
study is reported \emph{separately} throughout the paper because its
protocol differs from the synthesis-based protocol of \StudyReplication. Aggregating the three protocols would conflate three
fundamentally different sampling mechanisms.

\paragraph*{Generated-sampling protocol.}
The the performance--energy corpus artifact ships measured configuration tables
but no precomputed sampling traces. We therefore use a
\emph{generated-sampling} protocol: we draw samples from the measured
configuration universe using a common strategy panel and a budget
grid, then compute empirical $\TrainCost$ as the sum of the selected
configurations' measured property values. The protocol differs from
\StudyReplication (which reuse precomputed strategy traces from
Kaltenecker \etal) and from Linux (pool-based subsampling of
randconfig pools); the strategy set, in particular, is incompatible
with both.

\paragraph*{Selected systems.}
The full \StudyEnergy corpus covers $14$ systems; we report the
representative pass over $7$ systems chosen to span the cost-CV
range: \texttt{HSQLDB}, \texttt{brotli} (revision $1.0.7$),
\texttt{jump3r}, \texttt{lrzip}, \texttt{MongoDB}, \texttt{PostgreSQL}
(revision $11.2$), and \texttt{x264}. Pool sizes range from $180$
configurations (\texttt{brotli} 1.0.7) to $6{,}840$
(\texttt{MongoDB}), with no row caps applied. Note that the \StudyEnergy
\texttt{lrzip} and \texttt{x264} are \emph{not} the same artifacts as
the \StudyReplication \texttt{lrzip} or x264-matrix x264; we treat them as
separate observations under a separate protocol.

\paragraph*{Two cost properties, four scenarios per system.}
the performance--energy corpus measures both \emph{performance} and
\emph{energy} for each configuration. We use both as cost axes and
both as prediction targets, giving four scenarios per system:
$(\text{cost},\text{target})\in
\{\text{performance},\text{energy}\}^{2}$. RQ1 and RQ2 group by cost
property (they only ask about the cost axis); RQ3 and RQ4 are
reported per cost-target scenario.

\paragraph*{Strategies.}
Five deployable strategies plus two diagnostic oracles, all
pool-generated:
\begin{itemize}
    \item \texttt{random}: uniform random baseline.
    \item \texttt{balanced}: feature-balanced random sampling.
    \item \texttt{diversity}: cost-oblivious diversity in feature
    space.
    \item \texttt{pairwise}: greedy pairwise option-value coverage.
    \item \texttt{pilotCost}: pilot-based cost-aware diversity (uses
    a small pilot to estimate per-configuration cost, then samples
    diversely while accounting for cost).
    \item \texttt{cheapOracle}, \texttt{expensiveOracle}: lower- and
    upper-bound diagnostics that read the true cost signal.
\end{itemize}
None of the synthesis-based Kaltenecker strategies (Henard, t-wise,
distance-based, solver-based, diversified distance-based) is
evaluable here because there is no formal model handle, and only
\texttt{diversity} overlaps in name (but not in implementation) with
the Linux strategy of the same name. We therefore never combine.

\paragraph*{Sample sizes, learner, evaluation.}
A budget grid up to $n_{\max}=256$ with system-dependent
percentage anchors plus absolute anchors. Stochastic strategies are
repeated with $100$ seeds. The learner is ridge log-linear (same
family as \StudyReplication, Linux); MRE is computed on the complement
of the training sample within the measured universe.
We deliberately restrict the study to datasets in which
$\mathit{cost}(c)$ is a clean, well-defined wall-clock cost of running
configuration $c$ to obtain its label. Full-stack systems whose
acquisition cost is dominated by orchestration overhead (build, idle,
warm-up, monitoring) are out of scope of this paper because the cost
construct becomes the system's deployment workflow rather than the
sampling strategy's choice of configurations. ML and HPO benchmarks with
real runtime metadata (\eg LCBench, NAS-Bench-201) are an interesting
secondary direction but are deprioritized for this paper to keep the
empirical claim anchored in configurable software systems; we revisit
them in \Cref{sec:discussion} as future work. Surrogate-based ML
benchmarks (JAHS-Bench-201, YAHPO Gym, surrogate sub-benchmarks of
HPOBench) would weaken any empirical cost claim and are excluded.

\subsection{Sampling Strategies and Baselines}
\label{subsec:strategies}

We evaluate the six literature strategies introduced in
\Cref{subsec:sampling-strategies}: \emph{random},
\emph{solver-based}, \emph{randomized solver-based}, \emph{$t$-wise}
(coverage-based), \emph{distance-based}, and \emph{diversified
distance-based}. We use the implementations of Kaltenecker
\etal~\cite{kaltenecker2019}.

We also include the diagnostic baselines introduced in
\Cref{subsec:sampling-strategies}: a \emph{cheap-first oracle} and an
\emph{expensive-first oracle}, both selecting configurations in cost order
using known per-configuration times. These oracles are not deployable in
practice---they require pre-existing measurements---but they are
indispensable for diagnosing the cost-heterogeneity space: if cheap-first and
expensive-first behave very differently from random at the same $n$, there is
real cost heterogeneity to be exploited or stumbled into.

\subsection{Sample Sizes}
\label{subsec:sizes}

We evaluate strategies at a range of training sizes rather than at the natural
sizes of $t$-wise alone. For each system, we use percentage-based sizes
$\{1\%,2\%,5\%,10\%,20\%,30\%\}$ of $|C|$ where applicable, capped for very
small systems, and absolute sizes
$\{8,16,32,64,128,256\}$ as additional checkpoints. For $t$-wise we use
$t\in\{1,2,3\}$, record the resulting $n$, and include the corresponding
$(\text{strategy},n)$ pairs in the analysis. When all strategies can be
evaluated at a common $n$, we do so to enable direct comparison; otherwise we
analyze native strategy sizes separately.

\subsection{Learners}
\label{subsec:learners}

The primary learner is stepwise multiple linear regression, as in
Kaltenecker \etal~\cite{kaltenecker2019} and Pereira
\etal~\cite{pereira2020}. We use only configuration options as features---all
measured-property columns are excluded from features to prevent label
leakage. As a robustness check, we also report results with CART (consistent
with Sarkar \etal~\cite{sarkar2015}) and with a random-forest learner. The
primary results in \Cref{sec:results} use the literature-aligned learner; the
sensitivity to learner choice is reported in an appendix-style robustness
study so that the main story does not depend on a particular ML choice.

\subsection{Evaluation Protocol}
\label{subsec:protocol}

For each $(\text{system},\text{input},\text{property},\text{strategy},n,r)$
tuple, we
\begin{enumerate}
    \item generate the training sample $S_{r,s,n}$;
    \item compute $\TrainCost(r,s,n)$ from \Cref{eq:traincost};
    \item train the model on $S_{r,s,n}$ using configuration options as
    features only (all measured columns excluded; see below);
    \item evaluate on a fixed validation set $C_{\text{test}}$ or, when
    feasible, on $C\setminus S_{r,s,n}$;
    \item compute $\MRE$ and the count-based cost $\SizeCost=n$.
\end{enumerate}
We do \emph{not} include $C_{\text{test}}$ acquisition in the cost term: the
validation measurements are part of the experimental protocol and are the
same across strategies. They are not part of what the strategy controls.

\paragraph*{Forbidden columns.}
Strict feature hygiene is required because the cost we want to predict and the
cost we use to define $\TrainCost$ may overlap with the prediction target. For
each property under study, we exclude all measured-performance columns from
the feature set. For x264, when the target is \emph{time}, we exclude
\texttt{time}, \texttt{size}, and any related runtime/throughput columns; when
the target is \emph{size}, we exclude \texttt{size}, \texttt{usertime},
\texttt{systemtime}, \texttt{elapsedtime}, \texttt{cpu}, \texttt{frames},
\texttt{fps}, \texttt{kbs}, \texttt{etime}. Only configuration options enter
the features.

\paragraph*{Number of runs.}
For stochastic strategies (random, randomized solver-based, distance-based,
diversified distance-based), we use $100$ independent runs per
$(\text{system},\text{input},\text{property},\text{strategy},n)$ tuple,
matching the design of Kaltenecker \etal~\cite{kaltenecker2019} and
Pereira \etal~\cite{pereira2020}. Deterministic strategies (solver-based,
$t$-wise) are repeated only when seed or tie-breaking choices apply.

\paragraph*{Output schema.}
Each row of the raw results file records
\[
(\text{system},\text{input},\text{property},\text{strategy},n,r,
\TrainCost,\mathit{TrainCostPct},\mathit{SizePct},\MRE,\Ratio).
\]
Aggregations (medians, IQRs, statistical tests, correlations, $\lambda$-sweep
rankings, Pareto fronts) are computed from this raw table.

\paragraph*{Reproducibility.}
All code, raw measurement data, sampling implementations, and analysis
scripts will be made available on \red{anonymized URL}. We document seeds,
software versions, and the exact column-exclusion lists per system, so that
the analysis is end-to-end reproducible from raw measurements to plots and
tables.
